\section{Related work}
\label{sec:related}

Only works verified in \code{upstream-99p/docs/CITATION\_AUDIT.md} or
\code{docs/step0-trackA-findings.md} are cited. Several works that shaped the
methods are not yet verified and are left uncited; they are listed at the end of
this section by name only, and \code{refs.bib} marks them.

\paragraph{Benchmarks and models.}
Fail2Drive~\cite{gerstenecker2026fail2drive} pairs each unseen scenario with an
in-distribution twin on the same road. Its authors report SimLingo's sharpest
drop on PedestriansOnRoad and attribute it to the model following a
nonexistent car or cyclist (\code{docs/step0-trackA-findings.md}, A2). We use
that as motivation only: an authored scenario type is not an unlistable factor.
Bench2Drive~\cite{jia2024bench2drive} is a closed-loop driving benchmark on
CARLA~\cite{carla2017}, the simulator every pool here was driven in. SimLingo
\cite{Renz2025cvpr} and TransFuser++~\cite{jaeger2023hidden,zimmerlin2024hidden}
are the dev and first held-out models.

\paragraph{Finding systematic errors without labels.}
Slice discovery groups the inputs a model gets wrong so that a person can name
each group. Domino~\cite{eyuboglu2022domino} discovers systematic errors with
cross-modal embeddings, Spotlight~\cite{deon2022spotlight} is a general method
for discovering systematic errors in deep models, George~\cite{sohoni2020george}
targets hidden subclasses inside coarse classes, and
SliceLine~\cite{sagadeeva2021sliceline} finds slices for model debugging with
linear algebra. Our
setting differs in two ways. The unit is a closed-loop route, not a single
input, and the groups to recover are scenario types of a driving benchmark, not
planted image slices. We know of no published method tested on recovering known
failure groups of a closed-loop driving policy
(\code{docs/research/literature-review.md}).

\paragraph{Reading a policy's internals.}
Sparse autoencoders on vision-language-action policies give features that can be
inspected and, in some cases, steered~\cite{swann2026sparse,jin2026eventsae}
(\code{docs/step0-trackA-findings.md}). Dr.~VLA's per-feature
statistics, which separate general features from episode-specific ones, are the
basis of our B7 rival; using them as a gap signal is our construction. On driving
models, interpretability work has examined end-to-end models
\cite{motzkus2026driving}, probed and pruned the planning token of a driving
vision-language-action model~\cite{babu2026depthwise}, interpreted road scenes
through such models~\cite{panda2026whatdotheysee}, and extracted control vectors
from motion transformers~\cite{tas2025words}. Linear probes~\cite{alain2016probes}
underlie our probe diagnostics. None of these scores the recovery of known
failure groups.

\paragraph{Finding failures by search.}
Red-teaming and scenario generation look for new failures rather than group
known ones. Predictive red teaming~\cite{roboart2025} edits observations and
predicts degradation for manipulation policies. KING~\cite{hanselmann2022king},
AdvSim~\cite{wang2021advsim} and ChatScene~\cite{zhang2024chatscene} generate
safety-critical driving scenarios, CAT~\cite{zhang2023cat} trains against
adversarial scenarios in closed loop,
and a survey covers black-box safety validation~\cite{corso2021survey}. Failure
detectors such as SAFE~\cite{safe2025} predict that a failure will happen, which
is a different task from grouping failures by cause. Our methods could take
their input from any of these: once failures exist, discovery asks which belong
together.

\paragraph{Not yet verified, not cited.}
UniAD, the held-out model; caption-then-cluster failure discovery for robot
videos, the basis of H25 and the caption row; error-aware slice methods (FACTS,
B2T, InfEmbed, PlaneSpot); contrastive PCA (j02); cluster-wise stability (j08);
hidden Markov model sequence clustering (j09); Ward linkage; the bootstrap; and
DINOv2. Each needs a check against its primary source before it is cited.
